
\subsection{Is it possible for Norwegian Lutherans to build a Christian Free Church in America?}

Source: Kvartal-Skrift, 6th year (1880), pp. 180–190; 7th year (1881), pp. 23–33. — Ed.

\begin{quote}
\emph{Note: This section is currently represented by an editorial summary. 
A full translation of the original text is planned and contributions are welcome.}
\end{quote}


This section (pp. 101–119) asks whether Norwegian Lutherans can build a genuinely Christian free church in America, and it frames the question as both a response to criticism from Norway’s state-church press and a serious act of self-examination. The author notes that Norwegian commentators often mock or dismiss the Norwegian-American effort to gather around Word and Sacraments into free congregations and church bodies, while offering little understanding or constructive guidance. He finds it striking that, despite forty years of development and the existence of roughly a thousand Norwegian congregations and hundreds of pastors in America, leading figures in Norway’s state church show extraordinary indifference. The author suggests that this coldness cannot be explained merely by immigrant “strife” or incompetence; rather, it likely reflects fear and hostility toward what is most characteristic of the American Norwegian movement: “the free congregation in the free church.”

He argues that state-church leaders in Norway increasingly speak as though political and ecclesial freedom are inherently dangerous—treating republican government, church–state separation, congregational influence in calling pastors, lay testimony in public meetings, Sunday schools, and similar practices as products of unbelief or “free-thinking.” Faced with this, the author proposes that immigrants must ask honestly: does working for congregational freedom undermine Christian faith and lead to spiritual decay? If so, they would need to seek a state-supported church; if not, they must commit themselves more fully to rooting faith deeply through Word and Sacraments so it can withstand cultural winds.

He then examines America’s republican form of government and the separation of church and state. Republicanism, he contends, transfers responsibility from monarchs to the people, cultivating independence, accountability, and self-discipline—traits not opposed to Christianity’s call to awaken from spiritual sleep. The more contested issue is church–state separation, since many assume the church needs state support, forced revenues, and state-sponsored religious schooling. The author counters that it is neither unchristian nor harmful for Christians to bear the church’s costs themselves; indeed, it restores the church to a more apostolic pattern of self-support and voluntary sacrifice. He insists that coercive “state-church” means have not saved souls; true spiritual influence comes through God’s Word and Christ’s love. He appeals both to Christ’s own example—refusing political power and winning people through truth and mercy—and to American experience, where despite the absence of a state church the nation has been remarkably filled with churches and preaching.

The argument then turns to the deeper question: are free congregations themselves compatible with true Christianity? The author insists the decisive authority is the New Testament, not the practical “expediency” standards of state-church systems. While the New Testament is not a legal code of church constitutions, it clearly portrays the apostolic church as a fellowship of local congregations—free and self-governing yet united by the same Gospel, baptism, Lord’s Supper, Lord, and Spirit. Apostles are presented not as rulers but as servants; their authority is the authority of the Word, not a centralized church government. Each local congregation is fully “God’s congregation,” constituted by the right use of Word and Sacraments, and therefore possesses the same essential rights and responsibilities as the whole church.

Applying this to America, the author argues that immigrants were forced by circumstance to do what the apostolic church did: gather, call pastors, build churches, receive members, order worship, and carry responsibility themselves. This is not “humbug” but a natural and legitimate outworking of Christian freedom. Yet he warns of a new danger: a “dead free church” model associated with Missourian practice, where congregational freedom is reduced to the freedom to choose a pastor—and then effectively handed over to him. True congregational freedom, he argues, includes the flourishing of the Spirit’s gifts among all believers. Men and women should have room to serve in many ways—prayer, exhortation, mercy, teaching, practical administration—with appropriate offices like deacons and deaconesses where needed. Such participation is presented as the proper form of Christianity’s life in the world.

The section closes by stressing one more necessity for a durable free church: a serious pastoral seminary rooted in the free-congregation spirit and grounded in the catechetical “Børnelærdom.” With congregational freedom, catechesis, and a strong seminary, the author believes the Conference is on a sound path toward building a Christian free church in a free people.